\section{Experimental setup}
\label{sec:setup}

\subsection{Simulated acquisition and data}

We simulate linear fluorescence in a homogeneous, non-scattering medium. Numerical aperture (NA) describes the angular range admitted by the optics; we use $0.8$ for excitation and $0.1$ for detection. The low detection NA keeps a broad depth range similarly sharp, leaving weak depth cues in the camera image. For simplicity, we use light of wavelength $532$\,nm for both excitation and emission.

The phase plate has phase elements on a $0.32\um$ grid and lies $6.6\um$ before the intermediate illumination image plane; both values are expressed at specimen scale. We compare an untrained random plate with plates learned using our criterion, depth Fisher information, worst-case coherence or end-to-end training. The untrained plate has independent phases uniform on $[0,2\pi)$.

We average reference volumes into 16 depth planes of $48\times48$ voxels, each $0.64\um$ wide ($N=36{,}864$). The volume covers $30.72\times30.72\times10.24\um^3$. Each camera image supplies $12\times12$ samples. With $k$ illumination angles, measurement compression is $C=N/(144k)$: $256$, $64$ and $32$ images give $C=1$, $4$ and $8$. For each compressed setting, five random subsets of illumination angles are shared across plates. Finer simulation grids are specified in the supplement.

We normalize each plate's excitation to the same mean over the specimen and full angular scan. Scans with fewer images use proportionally longer exposures to match total illumination exposure, or nominal excitation dose, in expectation over random angle subsets. For each specimen, we calibrate the untrained full scan to a mean of 1000 expected photons per camera pixel per image and share this calibration across plates. Measurements include Poisson shot noise. Detected counts still depend on how illumination overlaps with the specimen, and peak intensities are not matched.

We simulate fine-scale neuronal volumes by rescaling public two-photon stacks from DEEP$^2$ \cite{wijethilake2023,wijethilake2023data}. Crops are normalized and resampled with axial extent reduced 3.125-fold; distances use simulation coordinates. The 60 test crops cover six cells from three animals. Each cell contributes two reference-selected and eight random crops, the latter sampled after reconstruction parameters were fixed. Other full-field studies use the 12 reference-selected crops unless stated otherwise.

Model-only optical training uses no specimen examples. Neuron-trained denoisers and supervised models use 459 training crops from a fourth animal. Supervised checkpoint selection uses 48 validation crops from a held-out stack of that animal. Two vascular crops from a separate animal are used to select TV/PnP scalar parameters. No test volume is used for training or tuning.

\subsection{Evaluation measures}

A depth profile records fluorescence intensity along depth at one lateral position. We report centroid error and one-dimensional Wasserstein-1 (W1) distance between normalized profiles \cite{peyre2019}. Both errors are in micrometers; lower is better.

For lateral column $c$, we normalize its reference and nonnegative reconstructed profiles to unit total intensity when nonzero. Let $F_c(i)$ and $\hat F_c(i)$ be their cumulative fluorescence fractions up to depth plane $i$. We compute
\begin{equation}
E_{\mathrm{W1}}=\sum_c\pi_c\,\Delta z\sum_{i=0}^{Z-1}\left|F_c(i)-\hat F_c(i)\right|,
\label{eq:w1}
\end{equation}
where $Z=16$ and $\Delta z=0.64\um$. We retain columns with total reference intensity above 10\,\% of the maximum and normalize these totals to sum to one, giving weights $\pi_c$. Centroid error uses the same columns and weights. TV/PnP outputs are clipped to $[0,1]$ before evaluation; feed-forward outputs remain raw except that W1 uses their nonnegative part. The supplement gives the zero-mass profile rule, foreground SSIM mask and signed-centroid convention.

Profile normalization removes total brightness, so we retain PSNR and foreground SSIM to assess intensity fidelity; higher values are better. W1 does not measure lateral displacement or guarantee filament separation. We test separation directly with controlled filament crossings (\cref{sec:res:structure}).

For neuronal comparisons, we use the same crops and angle subsets at matched nominal dose under the same noise model. For compressed scans, we average scores over the five subsets within each crop. We estimate 95\,\% confidence intervals by bootstrap resampling the six cells, keeping each cell's crops together. Dependence between cells from the same animal limits population-level interpretation.
